
\documentclass[11pt, a4paper]{article}

% ===== ESSENTIAL PACKAGES =====
\usepackage{geometry}
\geometry{a4paper, margin=1.2in}
\usepackage{setspace}
\usepackage{array}
\usepackage{xcolor}
\onehalfspacing
\setlength{\parskip}{1em}
\usepackage{multicol}

\usepackage{catchfile}

% ===== WORD COUNT COMMAND =====
\newcommand{\wordcount}{%
    \ifnum\pdf@shellescape=1%
        \immediate\write18{texcount -1 -sum Oblig1.tex > wordcount.txt}%
        \CatchFileDef{\words}{wordcount.txt}{}%
        \words%
    \else%
        [Enable -shell-escape]%
    \fi%
}
\usepackage{mdframed}

\newmdenv[
    leftmargin=-50pt,
    rightmargin=0pt,
    innerleftmargin=3pt,
    innerrightmargin=0pt,
    innertopmargin=5pt,
    innerbottommargin=5pt,
    skipabove=5pt,
    skipbelow=5pt,
    linewidth=1pt,
    linecolor=black,
    topline=true,
    rightline=false,
    bottomline=false
]{sectionboxinner}

\newenvironment{sectionbox}[1][]{%
    \begin{sectionboxinner}
    \textbf{\large #1}
    \vspace{1pt}
    \newline
}{%
    \end{sectionboxinner}
}

\newmdenv[
    leftmargin=0pt,
    rightmargin=0pt,
    innerleftmargin=0pt,
    innerrightmargin=0pt,
    innertopmargin=3pt,
    innerbottommargin=3pt,
    skipabove=8pt,
    skipbelow=8pt,
    linewidth=1.2pt,
    linecolor=gray!60,  % Lighter gray color
    topline=false,
    rightline=false,
    bottomline=false
]{subsectionboxinner}

\newenvironment{subsectionbox}[1][]{%
    \begin{subsectionboxinner}
    \textbf{#1}
    \vspace{3pt}
    \newline
}{%
    \end{subsectionboxinner}
}

% ===== CUSTOM ENVIRONMENTS FOR PHILOSOPHICAL ARGUMENTS =====
\newenvironment{argument}{\begin{quote}\itshape\textbf{Argument: }}{\end{quote}}
\newenvironment{objection}{\begin{quote}\itshape\textbf{Innvending: }}{\end{quote}}
\newenvironment{reply}{\begin{quote}\itshape\textbf{Reply: }}{\end{quote}}

% ===== DOCUMENT BEGIN =====
\title{Kap. 17 - Den sosiale kontrakten - hvorfor trenger vi stat og myndighet?}
\author{Oliver Ekeberg}
\date{\today}

\begin{document}
\maketitle

\tableofcontents

For å være en statsborger, er vi bundet av en rekke krav. Blant annet å betale skatt, ikke stjele andres eiendom, kjøpe billett på trikken. Disse normative kravene binder oss like mye som de begrenser oss. Du vil også bli straffet om du ikke handler i tråd med dem også.
\begin{itemize}
    \item "Du utviser ikke dyden pålitelighet" ville dydsetikeren sagt
    \item "Du respekterer ikke den andre som mål i seg selv" ville kantianeren sagt
    \item "Du setter størst mulig lykke prinsippet til side for egeninteresse" ville utilitaristen sagt
\end{itemize}

\begin{sectionbox}[Mennesket i naturtilstanden]

\begin{flushleft}
    \textbf{\textit{Staten har makt}} er grunnen til at vi aksepterer en slik inngripen at vi lar oss arrestere når vi bryter en lov. Denne makten kan staten bruke til å fremme hele statens interesser. Da kan vi spørre oss selv om det er en \textit{legitim} makt. Dette betyr at staten har en rett til å bestemme over folk og styre.\\
    
\end{flushleft}

\begin{flushleft}
    \textbf{\textit{Kontraktsteori}} er et forsøk på å forklare hvordan politiske lover og regler oppsto og hva rettferdighet og legitim politisk autoritet består i. 
\end{flushleft}

\begin{flushleft}
    Hvordan var livet i naturtilstanden? Det var umoralsk, dominert av egeninteresse, og du står i risiko for å miste huset, eiendeler eller livet. Dette er ifølge Thomas Hobbes, en filosof fra 1600 tallet under borgerkrigen i Storbrittania. Under naturtilstanden, vil folk søke seg sammen for å ikke utgjøre et offer for umoralske vesener. Ifølge Hobbes så. blir moralen, altså hva som er rett og galt til i samfunnet. Han ser for seg at det. er bedre i et samfunn, uansett styreform, enn naturtilstanden.
\end{flushleft}


\begin{flushleft}
    John Locke har et. mer optimistisk syn enn Hobbes.  Han ser for seg en naturtilstand hvor man har rettigheter, frihet og hvor moral har gyldighet. Mennesket er moralsk uten et samfunn fordi det er skapt av gud. Den viktigste grunnen til at vi trenger en stat, er at staten skal sikre våre naturlige rettigheter og plikter håndheves. Dette er fordi det er vanskelig å å følge og vurdere plikter og rettigheter når man selv er en part. Omformulert, det som oppnås når vi forlater naturtilstanden til fordel for politisk tilstand, er at rettferdighet blir sikret ved at dn blir institusjonalisert. Dvs. at en rettsstat gir oss stabile, upartiske institutsjoner som sikrer de rettighetene vi alle har. Staten er legitim om den er i tråd med sitt overordnede formål: å være til beste for alle
\end{flushleft}

\end{sectionbox}


\begin{sectionbox}[Eiendomsretten og penger]
\begin{flushleft}
John Locke mente at eiendomsretten er hovedmålet med det politiske samfunnet. Eiendomsretten inkluderer også retten på frihet og liv. Den viktigste eiendommen er oss selv. \\

Når blir noe mitt i naturtilstand? Når jeg sliter, dyrker jorden og blander mitt slit med omgivelsene. Det er bare min egen innsats som rettferdiggjør at jeg erklærer noe for mitt. Mine behov er også en \textit{naturlig begrensning på eiendomsretten}. Fordi jeg kan ikke spise mer enn jeg orker, derfor er det ikke noe vits i å dyrke mer enn jeg selv spiser. Det som forandret alt er penger. Penger har ingen holdbarhetsdato. Dermed gjør penger at vi får press på våre felles ressurser. Hvis grønnsaker kan byttes mot penger, har jeg insentiv til å dyrke mere med jorden min.\\

\textit{Kombinasjonen av eiendomsretten som en grunnleggende rettighet og kontrakten som grunnlaget for en legitim stat, gir oss byggesteiner for den teoretiske forståelsen av moderne liberale stater.} Bland inn penger, så muliggjør du fremvektsten av moderne kapitalisme. Slik sett er Lockes tankegods ekstremt viktig for liberal politisk teori og selvforståelsen til demokratiske kapitalistiske samfunn.
    
\end{flushleft}
\end{sectionbox}

\begin{sectionbox}[Samtykke og kontrakten]


Locke mener at vi automatisk gir samtykke til sosialkontrakten ved å oppholde oss i et stats territorium.
\end{sectionbox}

\end{document}